% ~1.25 pages. Ported 2026-11-21 from ../sec_methodology.tex (cgan/ run003 method) + ../sec_setup_operating_point.tex
% (operating point, environments, checked in full text in Session A, 2026-09-30) and updated to final/:
% roster = final/evaluate.py; recipe = final/train_gan.py; CNN = final/train_cnn.py; metrics = README Sec. 2.7;
% matched damage = final/figures.py matched(); gates = README Sec. 3.3q. Values belong here, not in Sec. system.
\section{Evaluation Methodology}\label{sec:method}

\subsection{Attackers}\label{sec:method:attackers}
Table~\ref{tab:attackers} lists the attackers, ordered by the information they use.
Three come from Li et al.'s jammer types~\cite{li2022jamming_ofdm}: barrage noise, a tone (a constant complex baseband value, its phase drawn per frame) and a pulse comb (an impulse every 64 samples).
The pulsed jammer transmits the victim's modulation and pulse with its own random symbols, each on with probability $p$ at amplitude $1/\sqrt p$, so that its mean power does not depend on $p$.
At $p=1$ it is the matched jammer that Zhou et al.\ call optimal~\cite{11184988}, at $p=0.25$ it is Li et al.'s protocol-aware type, and $p<1$ spans the pulsed family that contains the error-maximizing jammer against known modulation~\cite{amuru2015optimal}.
Its pulses ignore the frame, so at $p=0.01$ about $28\%$ of the frames receive none.
The frame-timed pulsed jammer uses the frame timing that the generator exploits (Sec.~\ref{sec:attacker}): it sends exactly $k\in\{1,2,3\}$ random symbols at random positions in every frame, with the frame's whole budget on them.
% attacks.pulsed_sync_tx; evaluated only against the deployed classifiers' eval files (eval_cnn_s13/14.json).
The on/off jammer uses the matched jammer at $+15$\,dB on a fraction $f$ of the frames and stays silent on the rest; its damage and $P_{\mathrm{det}}$ follow from the $+15$\,dB point by averaging over frames.
% analytic: P_det = f P_det(+15) + (1-f) FAR, damage = f damage(+15) (final/figures.py onoff()).
We call the lowest $P_{\mathrm{det}}$ that barrage, tone, pulse comb, the pulsed and the on/off jammer reach at a given damage the classical envelope; the frame-timed pulsed jammer is reported beside it.

Two references are not realizable.
The genie flip sends $-2\rho\,s[k]$ on a fraction $\delta$ of the symbols, synchronously; it inverts them, leaves the received law unchanged and so is invisible to any detector, which makes it the upper bound on damage at any $P_{\mathrm{det}}$.
The random push needs the genie's timing and carrier phase but neither its data nor the channel gain: it adds $1.1\,u[k]$, with $u[k]$ a random QPSK symbol, to a fraction of the symbols.
% genie push (eta 0.1) is table-only; random push: +1.1 u on floor(delta N) symbols, delta = JSR/1.21 (README 3.3q #9).

Zhou et al.'s generator~\cite{11184988}, trained to imitate the victim's waveform without a link or detector, is evaluated as published, in our reproduction.
Our own generators are trained per environment: a damage-only control ($\beta=0$, three seeds) and detector-aware generators with $\beta=10$ against one-sided energy and against kurtosis (one seed each) and against the classifier (three seeds).
% seeds: README 3.4 decision 7. "What the detector term buys" = cnn_b10 - control, three seeds each.
The control shares everything with the detector-aware generators except the detection term, so their difference is what that term buys.

\begin{table}[t]
\caption{Attackers, ordered by the information they use}\label{tab:attackers}
\centering\footnotesize
\begin{tabular}{@{}lll@{}}
\toprule
Attacker & Information & Parameters \\
\midrule
Genie flip & T0: symbols, sync & fraction $\delta$ \\
Random push & timing and phase & fraction \\
Pulsed~\cite{amuru2015optimal} & T1 & $p\in\{1,0.5,0.25,0.1,0.05,0.02,0.01\}$ \\
Pulsed, frame-timed & T1 & $k\in\{1,2,3\}$ per frame \\
On/off & T1 & fraction $f$ of frames \\
Zhou et al.\ CGAN~\cite{11184988} & T1 & as published \\
Generators (ours) & T1, simulation & target, $\beta\in\{0,10\}$ \\
Barrage, tone, pulse comb~\cite{li2022jamming_ofdm} & T2 & -- \\
\bottomrule
\end{tabular}
\end{table}

\subsection{Training the Generators}\label{sec:method:training}
The generator takes its layout from Zhou et al.~\cite{11184988}, a convolutional network with a dense output layer of about $8.6\times10^6$ parameters in all, but not their training: it starts from random weights and learns only from a differentiable form of~\eqref{eq:problem}, back-propagated through the link and the frozen detector.
No discriminator is trained.
% Zhou's imitation start biases G toward the victim's modulation: CNN P_det 0.48 vs 0.15 at matched BER,
% 4 seeds, 30 dB (README 3.3f) -- thesis appendix.
Two quantities are replaced by surrogates.
First, $\mathbb{E}[\mathrm{BER}]$ is numerically zero wherever the jammer is weak, so we maximize $\log\mathbb{E}[\mathrm{BER}]$, computed exactly over the receiver noise: a bit whose noiseless matched-filter sample lies at signed distance $m_b$ from its decision boundary is in error with probability $Q(m_b/\sigma)$, with $\sigma^2=\nu N_0/2$, and the logarithm of the mean is evaluated with a log-sum-exp so that it never underflows.
Second, $P_{\mathrm{det}}$ is replaced by
\begin{equation}
\tilde P_{\mathrm{det}}=\mathbb{E}\Bigl[\bigl(1+e^{-(T(r)-\tau_\alpha)/s_T}\bigr)^{-1}\Bigr],
\label{eq:softpdet}
\end{equation}
with $s_T$ the standard deviation of $T$ on clean frames; it tends to the hard detection probability as $s_T\to0$.
The generator minimizes
\begin{equation}
\mathcal{L}(G)=-\log\mathbb{E}[\mathrm{BER}]+\beta\,\tilde P_{\mathrm{det}},
\label{eq:loss}
\end{equation}
so $\beta$ takes a different numerical range than in~\eqref{eq:problem}.
The power constraint~\eqref{eq:power} is applied inside the forward pass: it removes the generator's overall gain, so the gradient acts only on the shape of the waveform.
The classifier's input image maps the spectrogram onto a 256-level colour map, whose gradient is zero; a straight-through estimator~\cite{bengio2013estimating} keeps the exact image in the forward pass and substitutes the colour map's slope in the backward pass, so the attacker differentiates the deployed network itself.

At each step the received JSR is drawn uniformly in decibels from a band that starts below the JSR at which the target detector begins to flag in that environment, because where it flags every frame the logistic term in~\eqref{eq:softpdet} is flat and carries no gradient.
At the ideal link this band is $-39$ to $-7$\,dB for the classifier and $-25$ to $0$\,dB for energy; the control uses $-32$ to $0$\,dB.
% bands.py: the 30 dB band shifted by how far the target's white-noise transition moved (README 3.4 step 2).
We use Adam with learning rate $2\times10^{-4}$ for $4000$ steps of $128$ frames ($32$ for the classifier), and train on the environment's own fading and noise-level draws.
% \todo{per-environment bands for all six, from bands.json?}

\subsection{Operating Point and Environments}\label{sec:method:env}
All experiments run at $E_s/N_0=24$\,dB, with $N=128$ symbols per frame, $L=8$ samples per symbol and a roll-off of $0.35$.
We take the victim to be the command-and-control link of an unmanned aircraft, which uses one fixed modulation; prototype radios for it use Gaussian minimum-shift keying~\cite{bishop2017cnpc}, for which QPSK, with a closed-form error rate, stands in.
A line-of-sight budget for such a link closes at the edge of its range at an SNR of $14$\,dB~\cite{kakar2015uav}; our operating point lies $10$\,dB above that edge, where uncoded QPSK makes practically no errors without a jammer, so the errors we count are the jammer's.

Two impairments the defender cannot know are added, first alone and then together.
Air-to-ground measurements in the bands allocated to these links~\cite{kerczewski2015cosite} found Rician K-factors of $12$--$15$\,dB in L-band and $27$--$31$\,dB in C-band across terrain types~\cite{matolak2017airground1,sun2017airground2,matolak2017airground3,khawaja2019survey}; we use $K=12$\,dB, the strongest L-band fading, and $K=28$\,dB.
A receiver's noise level is uncertain to about $\pm1$\,dB after calibration~\cite{shellhammer2006noise,tandra2008snr}; we model it as log-normal~\cite{sonnenschein1992radiometric,he2017covert} with $\sigma_N=0.5$\,dB, which keeps $95\%$ of frames within $\pm1$\,dB~\cite{nikonowicz2019noise}, and use $\sigma_N=1$\,dB for a receiver facing in-band interference.
This gives six environments: the ideal link (\emph{base}), noise-level uncertainty at $0.5$ and $1$\,dB (\emph{noise}, \emph{noise\,1\,dB}), fading at $K=12$ and $28$\,dB (\emph{fading}, \emph{fading\,K28}), and both at their main levels (\emph{both}).
In each, every threshold is set from $20\,000$ clean frames of that environment at $\alpha=0.05$ per frame, the value at which learned jamming detectors have been evaluated~\cite{arcangeloni2024jamming}.
The classifier is retrained in each environment by Li et al.'s recipe~\cite{li2022jamming_ofdm}: EfficientNet-B0, two classes, $762$ clean and $204$ jammed images per jammer type with the JSR uniform over $-20$ to $+10$\,dB, $100$ epochs.
We train three classifiers per environment that differ only in their random seed.
The generators are trained against the first, which is the attacker's surrogate and also sets their JSR band, and every classifier result is the mean over the other two, the deployed classifiers; the surrogate gives the white-box reference.
% SGD lr 1e-3, batch 32, 70/30 split (final/train_cnn.py). Surrogate seed 11 (cnn/), deployed seeds 13, 14
% (cnn_s13/, cnn_s14/; jobs 2454497-2454521). README 3.3q "CNN seeds 13 and 14". Li readout (Sec. res-li): surrogate
% seed 12, deployed seed 11.

\subsection{Metrics}\label{sec:method:metrics}
Every attacker is swept over the JSR from $-50$ to $+15$\,dB in $1$\,dB steps, with $512$ frames per point.
Damage is the excess bit and frame error rate (BER, FER) over clean frames of the same environment.
Because the jammer chooses its power, we compare attackers at matched damage: $P_{\mathrm{det}}$ is read at the first JSR where the excess damage reaches a level, interpolated in log damage.
We report it at an excess FER of $0.1$ and an excess BER of $3\times10^{-4}$; at $P_{\mathrm{det}}=0.1$ the binomial standard error of a point is $1.3$ percentage points.
Seeded generators are reported as mean $\pm$ standard deviation over seeds.
The simulation is checked against closed forms before every run: the clean and white-noise error rates, the matched-filter energy, the fading and noise-level laws, and each detector's realized false-alarm rate on fresh clean frames, which lies between $0.043$ and $0.057$ in every environment.
% final/verify.py 100/100 + regression gate vs an earlier independent run: 0/2178 P_det pairs beyond the
% sqrt(2)-corrected 4 sigma band (job 2278412). One sentence max in the paper.

\ifthesis
Before this design, we tried to train jammers by policy-gradient reinforcement learning over raw I/Q samples against a frozen classifier, which never learned; Appendix~\ref{app:history} records that and every other step.
\else
Policy-gradient reinforcement learning over raw I/Q samples, tried first, never learned against a frozen classifier; direct gradient through the detector does.
\fi
% README A.5: every sampled action flagged -> zero advantage -> no gradient.
